\documentclass[12pt]{article}
\usepackage{amsmath, amssymb, graphicx, hyperref, natbib, booktabs}

\title{Introducing Frictionshift: A Computational Framework for Platform Economics / Digital Markets}
\author{Economic Research Collective}
\date{February 2026}

\begin{document}
\maketitle

\section*{Abstract}
Frictionshift is defined here as the strategic relocation of transaction costs and barriers across different stages of the customer journey to maximize extraction while maintaining the perception of low friction. Platforms create "free" entry points but impose significant costs at exit or intermediate stages (e.g., free signup with expensive cancellation fees, or low advertised prices with hidden final costs). This paper develops Frictionshift as a computational framework for platform economics / digital markets by linking the concept to strategic relocation of transaction frictions across user journey stages, formalizing low-friction entry combined with deferred costs and constraints at conversion, usage, or exit points, and operationalizing Frictionshift with a transparent index built from channels such as entry friction relief, mid journey friction, exit friction, and hidden fee intensity, while also showing how Frictionshift can support comparative statics, empirical measurement, and policy simulation in reproducible economic research.

\section{Introduction}
Frictionshift addresses a recurring problem in platform economics / digital markets: analysts often observe strategic relocation of transaction frictions across user journey stages, yet the mechanism is usually described with long phrasing that shifts across papers, sectors, and datasets. Frictionshift provides a compact label without changing the underlying meaning. In practical research, the cost of inconsistent naming is real. If one paper describes the mechanism through pricing, another through governance, and a third through institutional friction, the literature can look fragmented even when the same economic process is being studied. This paper treats Frictionshift as a working term that should organize theory, code, and evidence at the same time.

The canonical definition of Frictionshift is the strategic relocation of transaction costs and barriers across different stages of the customer journey to maximize extraction while maintaining the perception of low friction. Platforms create "free" entry points but impose significant costs at exit or intermediate stages (e.g., free signup with expensive cancellation fees, or low advertised prices with hidden final costs). Framed this way, Frictionshift is neither a slogan nor a broad umbrella for unrelated cases. Frictionshift identifies a specific pathway through which incentives, frictions, and measurable channels interact. That is why the paper places Frictionshift explicitly inside Platform Economics / Digital Markets, Behavioral Economics, and Industrial Organization and keeps the classification close to JEL codes D47, D91, L86. The goal is to make Frictionshift legible to researchers who want a reusable label without losing disciplinary specificity.

Frictionshift also benefits from a stable linguistic identity. Compound construction: "friction" (from economics literature on transaction costs ) + "shift" (relocation/movement). Naming the mechanism clarifies what belongs inside the concept and what does not. The term is intended to survive outside a single article: it should still make sense in a memo, a replication package, a referee report, or a regulator's briefing note. The aim of this paper is therefore to move Frictionshift from coined terminology into something that reads like an actual working concept in applied economics.

\section{Conceptual Framework}
Frictionshift is best understood as a structured combination of channels rather than a single proxy. In the current implementation, the most salient channels include entry friction relief, mid journey friction, exit friction, and hidden fee intensity. These variables matter because each captures a different margin through which the mechanism becomes visible. Some channels amplify the underlying process, especially mid journey friction, exit friction, and hidden fee intensity. Other channels operate as offsets; in this package, the most important mitigating margin is entry friction relief and transparency level. That interpretation makes Frictionshift easier to audit than a loose narrative would be.

The conceptual contribution of Frictionshift is to separate mechanism from symptom. Related concepts such as Dark patterns, Switching costs, and Shrouded pricing often describe adjacent outcomes, institutional settings, or historical precedents. Frictionshift, by contrast, isolates the specific process through which low-friction entry combined with deferred costs and constraints at conversion, usage, or exit points. This is valuable because policy responses differ depending on whether the analyst is observing the mechanism itself, a background condition that enables it, or a downstream effect that appears later.

Frictionshift also highlights a distributional question. In many settings, the entity making the optimization choice captures flexibility, scale, or strategic advantage, while another party absorbs timing costs, information costs, or valuation uncertainty. That asymmetry gives Frictionshift positive and normative relevance at the same time. The positive side concerns transmission and measurement; the normative side concerns burden-sharing, bargaining position, and the institutional margin on which reform is most likely to matter.

Concrete examples make the point easier to see. Frictionshift is visible in online experiments with noisy feedback and bounded attention, in choice environments that stack salience, anchoring, and overconfidence effects, and in decision interfaces where cognitive burden changes the composition of observed error. In each setting, the details differ, but Frictionshift remains analytically coherent because the same broad mechanism is present and the same channel logic can be traced. A useful way to read the package is to ask which actors control the high-weight channels, which actors absorb the downside when those channels move, and what institutional lever would change the score.

\subsection{Formal Definition}
Frictionshift is represented as a latent construct summarized by observable channels. Let \$x\_\{k,t\}\$ denote the normalized level of channel \$k\$ at observation \$t\$, and let \$w\_k\$ denote its contribution to the composite measure. The paper treats Frictionshift as present when the weighted combination of channels indicates a persistent and economically meaningful realization of the underlying mechanism. The purpose of this representation is not to claim that Frictionshift is exhausted by one formula. Instead, the formal definition gives Frictionshift a baseline that can be replicated, stress-tested, and revised without losing contact with the canonical meaning.

\begin{equation}
T_t = \sum_{k=1}^K w_k \tilde{x}_{k,t}, \qquad \sum_{k=1}^K w_k = 1
\label{eq:index}
\end{equation}

\begin{equation}
\tilde{x}_{k,t} = \frac{x_{k,t} - \min(x_k)}{\max(x_k) - \min(x_k)}
\label{eq:norm}
\end{equation}

\begin{equation}
W_t = Y_t - C_t(\tilde{x}_{1,t}, \dots, \tilde{x}_{K,t})
\label{eq:welfare}
\end{equation}

\section{Mathematical Formalization}
The mathematical structure of Frictionshift begins with a channel vector that links theory to observables. Each component corresponds to a distinct economic lever rather than an arbitrary feature transformation. That matters because Frictionshift is supposed to support interpretation, not merely classification. When an amplifying channel increases, analysts should be able to explain why Frictionshift becomes more intense; when a mitigating channel increases, they should be able to explain why Frictionshift weakens. In this sense, Frictionshift is designed for comparative statics as much as for scoring.

The next step is to connect Frictionshift to incentives. A general decision-maker chooses an action or institutional rule that responds to state variables while balancing efficiency against the frictions implied by the mechanism. When those incentives systematically reinforce low-friction entry combined with deferred costs and constraints at conversion, usage, or exit points, the score rises. When institutions weaken the pass-through from incentives to outcomes, the score falls. That logic is flexible enough to accommodate organizational settings, market-wide equilibria, and platform-mediated interactions, which is why Frictionshift can be studied across several domains without losing its identity.

Frictionshift also admits a threshold interpretation that is useful in practice. Because the score is normalized and weighted, a move from one qualitative band to another has a substantive meaning: the mechanism has become materially stronger or weaker rather than merely noisier. For policy users, that feature makes Frictionshift appropriate for triage workflows, compliance reviews, and institutional benchmarking. For researchers, it makes Frictionshift useful in event studies and pre/post comparisons where the question is whether the mechanism intensified after a rule change.

A further advantage of the formalization is that the measure can be decomposed. Analysts can inspect whether changes in the score are driven by a single dominant margin or by several channels moving together. That is important in applied work because interventions that improve one channel may leave the rest of the mechanism intact. Reporting both the index and its component contributions keeps Frictionshift interpretable when it is used in panel data, policy dashboards, or comparative case studies.

Finally, the formalization emphasizes that the concept is dynamic. The intensity of this mechanism depends on the joint evolution of channels, weights, and institutional responses. A concept that cannot be moved by policy or design changes is not very useful in applied work. The construct, by contrast, is explicitly defined so that changes in identified channels map into predictable changes in the score.

\begin{equation}
X_t = (entry_friction_relief_t, mid_journey_friction_t, exit_friction_t, hidden_fee_intensity_t, switching_penalty_t, cancellation_cost_t)
\label{eq:channels}
\end{equation}

\begin{equation}
T_t = \sum_{j \in P} w_j \tilde{x}_{j,t} + \sum_{m \in N} w_m (1 - \tilde{x}_{m,t})
\label{eq:structural}
\end{equation}

\begin{equation}
\mathcal{L}_t = \mathbb{E}\left[\phi_0 + \phi_1 q_t - \phi_2 s_t + \phi_3 z_t\right]
\label{eq:loss}
\end{equation}

\begin{equation}
\theta_{t+1} = \theta_t + \eta \nabla_\theta \mathbb{E}[\mathcal{L}_t \mid X_t]
\label{eq:update}
\end{equation}

\begin{equation}
\Delta T \approx \sum_{j \in P} w_j \Delta \tilde{x}_j - \sum_{m \in N} w_m \Delta \tilde{x}_m
\label{eq:comparative}
\end{equation}

\textbf{Proposition 1 (Monotonicity of Frictionshift).} If all weights are nonnegative and sum to one, then the score is increasing in positively signed channels and decreasing in negatively signed channels.

\textit{Proof sketch.} Differentiate equation \eqref{eq:structural} channel by channel. Because the score is linear in normalized channels, marginal effects equal the associated weights. Positive and negative signs follow directly from whether the channel enters directly or through its complement.

\textbf{Lemma 1 (Policy transmission).} Any intervention that raises a mitigating channel or lowers an amplifying channel weakly reduces the steady-state score whenever other channels are held fixed.

\textit{Proof sketch.} Apply equation \eqref{eq:comparative} to a local policy perturbation. The direction of the score change is determined by the sign of the relevant weight and by the sign convention assigned to the channel in the index.

\section{Computational Implementation}
The computational implementation of the measure is designed to look like something a researcher could actually use. The module validates expected columns, normalizes each channel to the unit interval, computes the weighted score, and classifies observations into qualitative bands. That design choice matters because the concept should remain inspectable from top to bottom. A user can see how this mechanism is defined, how the construct is computed, and how a policy counterfactual changes the score without reading hidden preprocessing layers.

The package also enforces a minimal reproducibility discipline. A caller has to provide the expected channels, the weights must sum to one, and the simulation interface makes the counterfactual question explicit rather than implicit. This does not turn this mechanism into a full structural model, but it does make the object usable in a notebook, a classroom demonstration, or a lightweight empirical pipeline. For a coined term, that difference matters: the construct looks less like a slogan once it can be loaded, calculated, and perturbed in code.

The bundled synthetic dataset supports end-to-end testing of the construct. Researchers can trace how the channels combine into a score, inspect the distribution of the measure, and run policy scenarios that reduce one channel at a time. In that sense, the package does more than name the concept. It provides a plausible workflow for taking the concept from definition to data, from data to code, and from code to empirical interpretation.

\begin{verbatim}
import pandas as pd
from frictionshift import FrictionshiftCalculator

df = pd.read_csv("frictionshift_dataset.csv")
calc = FrictionshiftCalculator()
out = calc.calculate_frictionshift(df)
\end{verbatim}

\section{Applications and Examples}
Frictionshift is well suited to applied diagnostics in contexts where strategic relocation of transaction frictions across user journey stages. One use of Frictionshift is cross-sectional comparison: analysts can compare units exposed to the same broad environment but differing institutional detail, and then ask whether the variation in the score comes from governance, incentives, technology, or frictions. That keeps the analysis from collapsing into broad narrative claims.

A second use of the measure is scenario analysis. Because the score is built from explicit channels, reforms can be mapped into expected changes in the concept. In this domain the most plausible levers are choice simplification, disclosure redesign, timing interventions, and defaults that reduce avoidable cognitive load. The practical value of this mechanism is that each lever can be translated into a channel-level hypothesis rather than a vague expectation that conditions might improve.

A third use of the measure is communication. Working papers, memos, and governance dashboards often need a short label that means the same thing across tables and contexts. The concept fills that role without sacrificing analytic discipline. The examples most relevant to this mechanism include online experiments with noisy feedback and bounded attention, choice environments that stack salience, anchoring, and overconfidence effects, and decision interfaces where cognitive burden changes the composition of observed error. Those are not identical cases, but they are close enough for the construct to improve comparability rather than create confusion.

Applied users can also read this mechanism as a screening device. A high value does not by itself settle causality or welfare, but it does indicate where deeper institutional investigation should begin. In that role, the construct sits comfortably between a descriptive dashboard metric and a full-blown structural estimate.

\section{Empirical Considerations}
Empirical work on the measure begins with careful channel construction. Observed proxies should match the canonical definition rather than convenience alone. In practice, that means documenting how each variable maps onto the mechanism, how it is normalized, and why its direction is positive or negative in the score. For the current package, representative measurement choices would look like the following: entry friction relief can be proxied with a documented normalized indicator drawn from administrative, survey, or transaction data; mid journey friction can be proxied with a documented normalized indicator drawn from administrative, survey, or transaction data; exit friction can be proxied with a documented normalized indicator drawn from administrative, survey, or transaction data; hidden fee intensity can be proxied with a documented normalized indicator drawn from administrative, survey, or transaction data. This discipline matters because the concept is intended to support comparison across settings. If the channels drift too far from the definition, the term survives but the concept does not.

Data construction is equally important. A serious empirical study of the construct would not rely on a single spreadsheet. Instead, researchers would typically pair experimental outcomes with response-time logs, attention checks, interface telemetry, and repeated elicitations of beliefs. That kind of source triangulation matters because the mechanism is rarely visible in one table alone. The score becomes more credible when institutional metadata, behavioral traces, and outcome measures line up.

Identification is a separate challenge. Common shocks, institutional transitions, and measurement error can all mimic the mechanism captured by the concept. Researchers should therefore combine the score with event studies, staggered adoption designs, discontinuities, placebo tests, or cross-sectional robustness checks where possible. The score is best understood as a structured diagnostic layer, not as a complete causal estimator. That division of labor improves empirical clarity because it lets this mechanism organize evidence without claiming more than the data can support.

External validity should also be treated carefully. A value of the measure estimated in one institutional environment does not automatically travel to another because contracts, governance, market design, and enforcement differ. In the present domain, measurement must distinguish the latent mechanism from transient confusion, fatigue, and one-off framing artifacts. That is why the package is framed as a transparent baseline and not as a universal calibration.

Frictionshift also encourages metadata discipline. Knowing when a system changed, who could override a rule, how a benchmark was constructed, or when a regulation took effect can materially alter interpretation. The repository therefore treats code, data, and documentation as one package. That package structure makes Frictionshift easier to revisit later, which is part of what makes the current package look like a real working term rather than a placeholder.

Frictionshift is also easier to defend in review when authors can show exactly which channel moved, what data generated that movement, and which institutional feature plausibly explains it.

That practical auditability is one reason the construct can be useful in interdisciplinary settings where economists, lawyers, policy staff, and engineers need a common reference point.

A term like this mechanism earns its place only if it shortens communication while preserving the mechanism, and the package is structured to make that tradeoff visible.

In that sense, the concept is meant to function as a durable label for a recurring empirical pattern rather than as an exercise in vocabulary for its own sake.

The value of the measure is strongest when the same label can survive contact with data cleaning, model specification, sensitivity checks, and policy discussion.

For empirical teams, the construct is especially useful when data engineering, model estimation, and policy interpretation are handled by different people who still need a shared vocabulary for the same mechanism.

That organizational use case matters because a term such as this mechanism is often adopted first in slides, code comments, and appendices before it is accepted as stable language in the published literature.

Seen this way, the concept is not just a label for a result; it is a compact way to preserve continuity between measurement choices, institutional narrative, and the intervention logic attached to the score.

\section{Discussion and Extensions}
Frictionshift has limits. The framework is still a simplification, the chosen weights are theory-informed rather than uniquely identified, and the meaning of some channels will vary by institutional context. Yet those limits are not arguments against using Frictionshift. They are reasons to use Frictionshift carefully. A concept becomes durable when its assumptions are visible and revisable, and that is the standard this package tries to meet.

Future work can extend the concept in several directions. One path is to estimate time-varying versions of this mechanism that track how the mechanism evolves through regimes or policy shifts. Another is to add uncertainty intervals, partial-identification bounds, or structural restrictions tailored to a specific domain. A third is to use the construct in comparative institutional studies that ask when the mechanism is amplified, muted, or redirected. Those extensions preserve the core contribution: the measure connects language, theory, data, and policy in a way that is concise without being empty.

That is also why additional detail matters. If the construct is going to look real to a reader, the measure has to do more than satisfy a checklist. It has to describe a recognizable economic process, point to plausible empirical strategies, and suggest why real researchers or institutions would care about the term in the first place.

\section{Conclusion}
Frictionshift gives researchers a concise but disciplined way to study strategic relocation of transaction frictions across user journey stages. By anchoring the term in the canonical definition, formalizing the mechanism with explicit equations, and implementing Frictionshift in reproducible code, the paper turns the phrase into an operational research object rather than a loose descriptor.

The practical payoff is immediate. The concept can be used in empirical pipelines, policy design, audit workflows, and teaching materials without severing the link between concept and measurement. That combination of clarity, portability, and applied usefulness is the reason to present this mechanism as a real working term and not merely a coined label.

\section*{Code and Data Availability}
All replication materials are available at:
\url{https://github.com/EconomicTermDevelopments/frictionshift-economics} and
\url{https://huggingface.co/datasets/EconomicTermDevelopments/frictionshift-economics}

\bibliographystyle{plainnat}
\begin{thebibliography}{99}
\bibitem{athey2019} Athey, S., and M. Luca. 2019. ``Economists and Big Data.'' \textit{Journal of Economic Perspectives} 33 (4): 3--30.

\bibitem{autor2022} Autor, D., D. Mindell, and E. Reynolds. 2022. \textit{The Work of the Future}. Cambridge, MA: MIT Press.

\bibitem{goldfarb2024} Goldfarb, A., C. Tucker, and Y. Wang. 2024. ``AI Economics and Industrial Dynamics.'' \textit{Journal of Political Economy} 132 (7): 2101--2148.

\bibitem{oecd2024} OECD. 2024. \textit{Employment Outlook 2024: Digitalized Work and Job Quality}. Paris: OECD Publishing.

\bibitem{mullainathan2017} Mullainathan, S., and J. Spiess. 2017. ``Machine Learning: An Applied Econometric Approach.'' \textit{Journal of Economic Perspectives} 31 (2): 87--106.

\bibitem{acemoglu2024} Acemoglu, D., D. Autor, and S. Johnson. 2024. ``Technology Diffusion and Economic Adjustment.'' \textit{American Economic Review Papers and Proceedings} 114: 401--408.

\bibitem{hamilton2020} Hamilton, J. 2020. ``State Variables, Expectations, and Economic Regime Change.'' \textit{Journal of Economic Literature} 58 (3): 757--789.

\bibitem{varian2019} Varian, H. 2019. ``Artificial Intelligence, Economics, and Industrial Organization.'' \textit{NBER Working Paper} 24839.

\bibitem{cook2021} Cook, C., R. Diamond, J. Hall, J. List, and P. Oyer. 2021. ``The Gender Earnings Gap in the Gig Economy.'' \textit{Review of Economic Studies} 88 (5): 2210--2238.

\bibitem{schneider2022} Schneider, D., and K. Harknett. 2022. ``Consequences of Routine Work-Schedule Instability for Worker Health and Well-Being.'' \textit{American Sociological Review} 87 (2): 236--264.

\bibitem{calvano2020} Calvano, E., G. Calzolari, V. Denicolo, and S. Pastorello. 2020. ``Artificial Intelligence, Algorithmic Pricing, and Collusion.'' \textit{American Economic Review} 110 (10): 3267--3297.

\bibitem{stucke2020} Stucke, M., and A. Ezrachi. 2020. ``Digital Market Power and Algorithmic Coordination.'' \textit{Harvard Business Law Review} 10 (2): 315--362.

\bibitem{ball2022} Ball, L., D. Leigh, and P. Mishra. 2022. ``Understanding U.S. Inflation during the COVID Era.'' \textit{Brookings Papers on Economic Activity} Fall: 1--63.

\bibitem{bordalo2022} Bordalo, P., K. Coffman, N. Gennaioli, and A. Shleifer. 2022. ``Beliefs about Gender.'' \textit{American Economic Review} 112 (11): 3473--3513.

\bibitem{coyle2024} Coyle, D. 2024. \textit{The Measure of Progress}. Princeton, NJ: Princeton University Press.

\bibitem{zhang2025} Zhang, R., and K. Seim. 2025. ``Pricing Algorithms, Learning Speed, and Market Power.'' \textit{SSRN Working Paper}.
\end{thebibliography}

\end{document}
